\documentclass[10pt,a4paper]{article}

\usepackage[top=0.5in,bottom=0.5in,left=0.62in,right=0.62in]{geometry}
\usepackage{enumitem}
\usepackage{titlesec}
\usepackage{hyperref}
\usepackage[T1]{fontenc}
\usepackage[utf8]{inputenc}
\usepackage[brazilian]{babel}
\usepackage{xcolor}
\usepackage{parskip}

\definecolor{accent}{HTML}{1A1A1A}
\definecolor{link}{HTML}{0056B3}
\hypersetup{colorlinks=true,urlcolor=link,pdfborder={0 0 0}}
\pagestyle{empty}
\titleformat{\section}{\normalsize\bfseries\color{accent}}{}{0em}{}[\titlerule]
\titlespacing*{\section}{0pt}{7pt}{3pt}
\setlist[itemize]{leftmargin=14pt,itemsep=1.5pt,topsep=2pt,parsep=0pt}

\newcommand{\job}[4]{%
  \textbf{#1} \hfill \textbf{#2}\\
  \textit{#3} \hfill \textit{#4}\\[-2pt]
}

\begin{document}

\begin{center}
{\Large\bfseries Yago Henrique Teixeira}\\[3pt]
\textbf{Desenvolvedor iOS | Swift | SwiftUI | Arquitetura Mobile}\\[3pt]
Belo Horizonte, MG \quad | \quad (31) 99834-8635 \quad | \quad
\href{mailto:yago.ht@icloud.com}{yago.ht@icloud.com} \quad | \quad
\href{https://www.linkedin.com/in/yagohtx/}{linkedin.com/in/yagohtx}
\end{center}

\section{Resumo Profissional}

Desenvolvedor iOS com 6 anos de experiência em desenvolvimento de software, especializado em Swift, SwiftUI e arquitetura de aplicações escaláveis. Atuo na evolução de um aplicativo B2B de e-commerce utilizado em múltiplos países da AB InBev, liderando iniciativas de modernização de código legado, arquitetura modular, testes automatizados, observabilidade e automação do fluxo de desenvolvimento com IA.

\section{Competências Técnicas}

\textbf{iOS e arquitetura:} Swift, SwiftUI, UIKit, Combine, async/await, Core Data, Realm, SwiftData, TCA, MVVM, Clean Architecture, Coordinator, Swinject, ViewCode, Deep Links

\textbf{Qualidade:} XCTest, Quick, Nimble, SnapshotTesting, ViewInspector, testes unitários, testes de snapshot

\textbf{Entrega e observabilidade:} Firebase Remote Config, AppLanga, Segment, New Relic, Crashlytics, CocoaPods, Swift Package Manager (SPM), Git, Azure Pipelines, Fastlane

\textbf{Backend e IA:} Cursor, agentes de IA, Skills, Model Context Protocol (MCP), Jira, Azure, Java, Spring Boot, Node.js, REST APIs, SQL, PostgreSQL, MongoDB, Apache Kafka, WebSockets

\section{Experiência Profissional}

\job{Desenvolvedor iOS}{Jun 2022 -- Atual}{AB InBev}{Belo Horizonte, MG}

\begin{itemize}

\item Desenvolvo e evoluo um aplicativo iOS B2B de e-commerce utilizado em múltiplos países, conectando varejistas ao catálogo de produtos da AB InBev.

\item Participei da adoção de um Back-end for Front-end (BFF), implementando componentes reutilizáveis em SwiftUI, filtros, trays e novas experiências de interface, além do mapeamento de respostas para a camada de apresentação e tratamento de erros com async/await.

\item Redesenhei a página de listagem de produtos utilizando The Composable Architecture (TCA), isolando estado, lógica de negócio e responsabilidades da interface.

\item Instrumentei eventos no Segment para ampliar a observabilidade da jornada do usuário, monitorando navegação, filtros, buscas e compras para apoiar decisões de produto.

\item Desenvolvo testes unitários e de snapshot para reduzir regressões, além de utilizar New Relic e Crashlytics na investigação de problemas de performance e falhas em produção.

\item Desenvolvo funcionalidades para múltiplos países utilizando localização, Firebase Remote Config, AppLanga, configurações específicas por mercado e testes A/B com Optimizely.

\item Atuo em um monorepo modular com nove módulos compartilhados entre diferentes equipes. Também implementei o Mockoon em um aplicativo interno de demonstração, reduzindo a dependência do backend durante testes locais.

\item Proponho soluções de arquitetura e automações com IA utilizando Cursor, Skills e MCP integrados ao Jira e Azure, acelerando refinamentos, implementação, documentação, versionamento e criação de Pull Requests.

\end{itemize}

\job{Desenvolvedor iOS Freelancer}{Jan 2024 -- Dez 2024}{Rhitmo Tech}{Belo Horizonte, MG}

\begin{itemize}

\item Liderei a modernização completa de um aplicativo iOS de controle de ponto, migrando uma base monolítica em UIKit, sem arquitetura definida, para MVVM com SwiftUI.

\item Migrei todas as telas do aplicativo para SwiftUI, preservando as regras de negócio existentes e reduzindo significativamente a complexidade de manutenção da aplicação.

\item Estruturei um sistema de componentes reutilizáveis inspirado em Atomic Design, padronizando interfaces e reduzindo duplicação de código.

\item Reestruturei a camada de networking, substituindo uma implementação centralizada composta por métodos extensos por uma arquitetura modular com responsabilidades bem definidas.

\item Reorganizei a persistência utilizando Core Data, desacoplei navegação, regras de negócio e interface e implementei testes unitários para as principais funcionalidades.

\item Reconstruí regras de negócio por meio de engenharia reversa das APIs e documentação incompleta, transformando um sistema legado sem arquitetura em uma base sustentável.

\item Atuei como referência técnica do projeto, mentorando dois desenvolvedores, participando da contratação técnica da equipe e conduzindo decisões de arquitetura.

\end{itemize}

\job{Desenvolvedor Java}{Nov 2020 -- Jun 2022}{Avenue Code}{Belo Horizonte, MG}

\begin{itemize}

\item Desenvolvi aplicações web utilizando Java e Spring Boot, implementando APIs REST, regras de negócio e integrações com SQL e MongoDB.

\item Trabalhei em fluxos assíncronos utilizando Apache Kafka e desenvolvi suítes de testes automatizados com Cucumber (BDD).

\end{itemize}

\job{Estagiário de Desenvolvimento Java}{Mai 2020 -- Nov 2020}{Avenue Code}{Belo Horizonte, MG}

\begin{itemize}

\item Desenvolvi serviços backend em Java, Spring Boot e Node.js para aplicações internas.

\item Implementei comunicação assíncrona com Apache Kafka para distribuição de eventos em uma plataforma social corporativa.

\end{itemize}

\section{Formação Acadêmica}

\textbf{Bacharelado em Ciência da Computação} \hfill 2017 -- 2020

Centro Universitário de Belo Horizonte (UniBH) \hfill CR: 91,29/100

\section{Informações Adicionais}

\textbf{Idiomas:} Português nativo; inglês fluente, utilizado diariamente em reuniões, documentação e colaboração com equipes internacionais.

\textbf{Práticas de trabalho:} Scrum, métodos ágeis, code review, pair programming, refinamentos técnicos e Planning Poker.

\end{document}